\capitulo{6}{Trabajos relacionados}

\section{Introducción}

Los sistemas SIEM son ampliamente utilizados en entornos empresariales, industriales y gubernamentales para la monitorización centralizada de eventos de seguridad. Existen numerosas soluciones comerciales y de código abierto que ofrecen funcionalidades avanzadas de recopilación de logs, correlación de eventos, detección de amenazas y generación de alertas.

A continuación se presentan algunas de las soluciones más conocidas relacionadas con el ámbito de los sistemas SIEM y se realiza una breve comparación con el sistema desarrollado en este proyecto.

\section{Splunk}

Splunk es una de las plataformas SIEM más utilizadas a nivel empresarial. Permite la recopilación y análisis de grandes volúmenes de datos procedentes de múltiples fuentes, proporcionando capacidades avanzadas de búsqueda, generación de informes y visualización.

Entre sus principales ventajas destacan su escalabilidad y la gran cantidad de funcionalidades disponibles. Sin embargo, se trata de una solución compleja y con elevados costes de licencia para determinados entornos.

\section{IBM QRadar}

IBM QRadar es una plataforma SIEM orientada principalmente a grandes organizaciones. Dispone de mecanismos avanzados de correlación, análisis de comportamiento y detección de amenazas.

Su principal ventaja reside en la capacidad para gestionar grandes infraestructuras y ofrecer análisis avanzados de seguridad. Como contrapartida, requiere una infraestructura considerable y una configuración más compleja que la necesaria para proyectos de menor tamaño.

\section{Wazuh}

Wazuh es una plataforma de código abierto ampliamente utilizada en entornos de monitorización y ciberseguridad. Ofrece funcionalidades relacionadas con la recopilación de eventos, detección de intrusiones, monitorización de integridad y generación de alertas.

Al tratarse de una solución de código abierto, constituye una alternativa muy interesante para organizaciones que buscan reducir costes manteniendo un elevado nivel de funcionalidad.

\section{Elastic Stack}

Elastic Stack, también conocido como ELK Stack (Elasticsearch, Logstash y Kibana), es una plataforma muy utilizada para la gestión y visualización de grandes volúmenes de información.

La combinación de estas herramientas permite construir soluciones similares a un SIEM mediante la recopilación, almacenamiento, análisis y representación gráfica de eventos.

\section{Comparación con el sistema desarrollado}

A diferencia de las soluciones anteriores, el objetivo de este proyecto no ha sido competir con plataformas profesionales ampliamente consolidadas, sino desarrollar un sistema SIEM ligero con fines académicos que permita comprender los principios fundamentales de funcionamiento de este tipo de herramientas.

El sistema desarrollado incorpora funcionalidades básicas de ingesta de eventos, almacenamiento, correlación, detección de anomalías, API REST y dashboard de monitorización. Aunque su alcance es más reducido que el de las plataformas comerciales, permite reproducir gran parte del flujo de trabajo presente en sistemas SIEM reales.

Además, la arquitectura modular utilizada facilita futuras ampliaciones e incorpora muchos de los conceptos fundamentales empleados por las soluciones profesionales analizadas.


\begin{table}[htbp]
\centering
\begin{tabular}{|l|l|c|c|c|}
\hline
\textbf{Solución} & \textbf{Tipo} & \textbf{Código abierto} & \textbf{Correlación} & \textbf{Dashboard} \\
\hline
Splunk & Comercial & No & Sí & Sí \\
IBM QRadar & Comercial & No & Sí & Sí \\
Wazuh & Open Source & Sí & Sí & Sí \\
Elastic Stack & Open Source & Sí & Parcial & Sí \\
TFG-SIEM & Académico & Sí & Sí & Sí \\
\hline
\end{tabular}
\caption{Comparativa entre diferentes soluciones SIEM}
\label{tab:comparativa_siem}
\end{table}

Como puede observarse, las plataformas comerciales ofrecen un mayor número de funcionalidades y están preparadas para gestionar infraestructuras de gran tamaño. Sin embargo, el sistema desarrollado en este proyecto incorpora los elementos fundamentales presentes en un SIEM moderno, permitiendo comprender de forma práctica los procesos de ingesta, almacenamiento, correlación y monitorización de eventos de seguridad.

